\documentclass[10pt]{article}

\usepackage[utf8]{inputenc}

\usepackage[T1]{fontenc}

\usepackage{amsmath}

\usepackage{amsfonts}

\usepackage{amssymb}

\usepackage[version=4]{mhchem}

\usepackage{stmaryrd}



\begin{document}

где сумма числа компонент векторов $y_{1}$ и $y_{2}$ равна $n$, существуют, вообще говоря, п. с. в окрестностях концов сегмента $[0, T]$. Построен алгоритм нахождения асимптотики решения этоӥ задачи, при определенных свойствах функцй $f$ и $g$ доказано существование и единственность решения и даны его оценки (см. [3]). В слутае неединственности решения предельного уравнения $g=0$ относительно $y$ построен внутренний (в окрестности $\tau, 0<\tau<T)$ п. с., разделяющий области с различными решениями предельного уравнения. Для одного типа интегро-дифференциального уравнения построен алгоритм асимптотич. разложения по $\varepsilon$ задачи с начальными условиями и исследованы нек-рые особенности поведения решений.



В случае линейных обыкновенных дифференциальных уравнений $(L+\varepsilon M) x=f(t), 0 \leqslant t \leqslant 1$, где $L$ и $M-$ дифференциальные операторы, с граничными условиями $B_{0}+B_{1}$ выделен класс задач, решение к-рых содержит п. с., и введено понятие регулярного вырождения, при к-ром решение предельного уравнения позволяет удовлетворить условиям $B_{0}$, а асимптотич. решение для п. с. - условиям $B_{1}$ (см. [4]). Построен итерационный процесс асимптотич. представления решения и даны оценки остаточных членов разложений.



В П. с. т. общего нелинейного обыкновенного дифференциального уравнения 2-го порядка при определенных предположениях доказано (см. [5]), что решение 1-й краевой задачи складывается из внешнего решения, п. с. и остаточного члена, имеющего с 1-й производной порядок $\varepsilon$ на всем сегменте.



Изучено поведение решений краевых задач основных типов для линейного уравнения с частными производными вида



\[

\varepsilon \Delta u+A(x, y) u_{x}+B(x, y) u_{y} \dashv C(x, y) u=f(x, y),

\]



где $\Delta$ - оператор Лапласа, в области $D$ с границей $S$. Построены условия на функции $A, B, C, f$, границу $S$ и функции $a, \varphi$ точки $P$ линии $S$, входящие в граничное условие $u_{n}+a(P)=\varphi(P)$, при к-рых $u$ в $D+S$ равномерно стремится к решению предельного уравнения с приведенным граничным условием на определенной части $S$ (отсутствие п. с.) (см. [6]).



Для эллиптич. уравнения 2-го порядка в области $D$ с границей $S$ на примере двух независимых переменных



\[

\begin{gathered}

\varepsilon\left[a(x, y) u_{x x}+2 b(x, y) u_{x y}+c(x, y) u_{y y}+\right. \\

\left.+d(x, y) u_{x}+e(x, y) u_{y}+g(x, y) u\right]+u_{x}- \\

-h(x, y) u=f(x, y), h \geqslant \alpha^{2}>0,

\end{gathered}

\]



построены итерационные процессы решения задачи с условием $u=0$ на $S$, доказаны теоремы о структуре разложения $u$ по $\varepsilon$ и даны оценки остаточного члена этого разложения (см. [4]). Аналогичные результаты получены и для уравнений высших порядков.



Разработан (см. [7]) метод сращивания асимптотич. разложений для уравнения



\[

\varepsilon \Delta u-a(x, y) u_{y}=f(x, y)

\]



в прямоугольнике с заданным $u$ на его границе.



Исследования П. с. т. для нелинсйных уравнений с частными производными связаны в основном с аэрогидродинамикой и базируются на уравнениях НавьеСтокса и их обобщениях. Запросы практики привели как к развитию математич. теории, так и к возникновению различных задач и методов их решения. Здесь речь будет идти только о ламинарных течениях (см. [8] - [10]).



Гидродинамика плоских $(k=0)$ и осесимметричных $(k=1)$ тетений несжимаемой жидкости с постоянным коэффициентом вязкости $v$ описывается уравнениями Навье - Стокса



\[

\left.\begin{array}{c}

\left(\eta^{k} u\right)_{\xi}+\left(\eta^{k} v\right)_{\eta}=0, \quad u_{\tau}+u u_{\xi}+v u_{\eta}=\varepsilon^{2} \Delta u-p_{\xi}, \\

v_{\tau}+u v_{\xi}+v v_{\eta}=\varepsilon^{2} \Delta v-p_{\eta},

\end{array}\right\}

\]



гпе $\varepsilon=R^{-1 / 2}, R=w X / v$ - число Рейнольдса, представленное через характерные величины скорости $w$ и линейного размера $X$. На замкнутой границе $S$ области $D$ решения задаются краевые условия, причем на обтекаемом контуре $\Gamma$ при $\eta=H(\xi)$ задаются условия $\bar{u}=0, \bar{v}=v_{0}(\xi, H(\xi))$, где $\bar{u}, \bar{v}-$ касательная и нормальная к $\Gamma$ составляющие вектора $(u, v)$. В $D+S$ задаются начальные значения $u, v, p$.



При малом $\varepsilon$ асимптотич. решение задачи в первом приближении составляется из решения уравнений (2) при $\varepsilon=0$ с частью условий на $S$ (на $\Gamma$ ставится только условие $v=v_{0}$ ) и решения уравнений п. с. Уравневия динамического п. с. выводятся в предположении, что условная толщина п. с. $\delta$ и величина $v$ имеют порядки $\delta \sim X \varepsilon, \quad v \sim w \varepsilon, \quad$ а члены левых частей последних уравнений из (2) имеют порядок членов с $\varepsilon^{2}$. Введение переменных $t=\tau, x=\xi, \quad y=\eta / \varepsilon, V=v / \varepsilon$ приводит при $\varepsilon \rightarrow 0$ к уравнениям Прандтля



$\left(r^{k} u\right)_{x}+\left(r^{k} V\right)_{y}=0, u_{t}+u u_{x}+V u_{y}=u_{y y}-p_{x}, p_{y}=0,(3)$ $0 \leqslant t \leqslant T, 0 \leqslant x \leqslant X_{0}, 0 \leqslant y<\infty$,



с условиями



\[

\begin{gathered}

\left.u\right|_{t=0}=u_{0}(x, y),\left.v\right|_{t=0}=v_{0}(x, y),\left.u\right|_{y=0}=0, \\

\left.V\right|_{y=0}=v_{0}(t, x),

\end{gathered}

\]



$u \rightarrow W(t, x)$ при $y \rightarrow \infty, p_{x}=-W W_{x}-W_{t},\left.u\right|_{x=0}=0$, где $r$ - расстояние от оси симметрии при $k=1, W(x)$ известная функция. Эти уравнения и условия справедливы для любого криволинейного контура, радиус кривизны к-рого много больше $\delta$. В последнем случае $x, y$ - координаты вдоль контура и по нормали к нему.



При постоянном $W$ задача сводится к краевой для обыкновенного дифференциального уравнения. Имеются и другие классы подобных решений.



Известны условия, при к-рых репения задач п. с. существуют, исследованы вопросы единственности и устойчивости решений и их выхода на решения стационарных задач (см. [11]). Решения строятся по методу прямых, показана их сходимость.



Уравнения п. с. для сжимаемой жидкости выводятся из уравнений для течевий вязкого и теплопроводного газа и значительно более сложны, чем (3). Возрастает и их количество. Имеется интегральное преобразование, упрощающее эти уравнения в общем случае и сводящее их к (3) при числе Прандтля $\operatorname{Pr}=c_{p} / K=1$, где $c_{p}$ - теплоемкость газа при постоянном давлении, $K$ - коэффициент теплопроводности (см. [12]). Известен ряд модификаций преобразования. В общем случае уравнения п. с. описывают т. н. естественные конвективные течения. Если $\boldsymbol{v}$ не зависит от температуры и архимедова сила пренебрежимо мала, то уравнекие әнергии отделяется от системы уравнений п. с. и говорят о вынужденных конвективных течениях. Уравнение энергии определяет тепловой п. с., толщина к-рого отличается от $\delta$.



П. с. возникают и в зоне раздела течений с различными характерными скоростями. Ударные волны также являются п. с.



Отдельный класс двумерных задач п. с. связан с течениями у вращающихся осесимметричных пластин и тел.



Вместе с развитием методов решения нестационарных задач решены задачи с периодическим $W$, при движении скачком из состояния покоя, при ускорен-





\end{document}